智慧农业研究正在从单一传感器观测转向多源数据协同分析。作物表型数据既包括田间图像、环境传感、人工调查记录，也包括由模型推断得到的生长指标。不同来源的数据在空间尺度、时间频率、噪声结构和解释方式上差异显著，因此需要在进入统计分析或机器学习模型之前完成一致化处理。

吉林农业大学研究生学位论文规范要求农学、理学、工学、医学类论文正文原则上包含文献综述、材料与方法、结果与分析、讨论、小结等部分。本示例按照该口径组织正文，并保留前言、结论、参考文献、附录、作者简介和致谢等后置结构，便于用户在真实写作时逐段替换。
